\documentclass[a4paper]{article}
\usepackage[margin=3cm]{geometry}
\usepackage{graphicx}

\usepackage{hyperref}
\hypersetup{
	colorlinks=true,
	linkcolor=blue,
	filecolor=magenta,
	urlcolor=cyan,
	hyperindex=true,
	linktocpage=true,%makes the page number be the link in the index
}
\usepackage{listings}
\lstset{
	basicstyle=\tiny\ttfamily,
	columns=fullflexible,
	keepspaces=true,
	breaklines=true,
	frame=single,
	showstringspaces=true,
	tabsize=4,
	showspaces=true,
	showtabs=true,
}

\renewcommand{\arraystretch}{1.2}

\title{Calutron refining}

\author{Chaya2766}

\date{last updated \today}

\begin{document}
	\sffamily
	\hyphenchar\font=-1
	\sloppy
	
	\maketitle
	
	\textbf{
		\sloppy
		\hyphenchar\font=-1
		}
	
	\vfill
	
	\begin{abstract}
		This document investigates refining of pure elements from general matter through a process involving the ionization of that matter and deflection of the ions in a magnetic field to separate them by mass.
		
		\medskip
		
		The approximate energy required to separate one kilogram of matter composed in equal parts by weight of the first 104 elements of the periodic table into the individual elements via this process is between 6.4685GJ and 64.685GJ.
		
		\medskip
		
		For meteor material that energy is between 15.5GJ and 155GJ.
	\end{abstract}
	
	\vfill
	
	\begin{quote}
		\begin{center}
			\textbf{DISCLAIMER AND WARNING}
		\end{center}
		
		This is not a scientific paper or engineering advice.
		This document is only intended as creative world‑building and has not been peer‑reviewed or approved for any practical application.
		Use of these ideas in real‑world designs or other serious application is entirely at the reader’s discretion and risk.
	\end{quote}
	
	%\noindent \rule{\linewidth}{1pt}
	\vfill
	
	\tableofcontents
	
	\pagebreak
	
	\section{Ionization energy}
	
	These values had been sourced from \href{https://ptable.com/#Properties/Ionization/1st}{Ptable}.
	
	\scriptsize
	
	\begin{center}
	\begin{tabular}{|c|c|c|c|c|}
		\hline
		atomic number & element name & molar mass (gram/mol) & 1st ionization energy (kJ/mol) & specific ionization energy (MJ/kg) \\
		\hline
		1 & hydrogen & 1.008 & 1312 & 1301.5873015873 \\
		\hline
		2 & helium & 4.002602 & 2372.3 & 592.689455509191 \\
		\hline
		3 & lithium & 6.94 & 520.2 & 74.9567723342939 \\
		\hline
		4 & beryllium & 9.0121831 & 899.5 & 99.809334765957 \\
		\hline
		5 & boron & 10.81 & 800.6 & 74.0610545790934 \\
		\hline
		6 & carbon & 12.011 & 1086.5 & 90.4587461493631 \\
		\hline
		7 & nitrogen & 14.007 & 1402.3 & 100.114228599986 \\
		\hline
		8 & oxygen & 15.999 & 1313.9 & 82.1238827426714 \\
		\hline
		9 & fluorine & 18.998403162 & 1681 & 88.4811205271337 \\
		\hline
		10 & neon & 20.1797 & 2080.7 & 103.108569503015 \\
		\hline
		11 & sodium & 22.98976928 & 495.8 & 21.5661146469757 \\
		\hline
		12 & magnesium & 24.305 & 737.7 & 30.351779469245 \\
		\hline
		13 & aluminium & 26.9815384 & 577.5 & 21.4035238257578 \\
		\hline
		14 & silicon & 28.085 & 786.5 & 28.0042727434574 \\
		\hline
		15 & phosphorus & 30.973761998 & 1011.8 & 32.6663580634904 \\
		\hline
		16 & sulfur & 32.06 & 999.6 & 31.17903930131 \\
		\hline
		17 & chlorine & 35.45 & 1251.2 & 35.2947813822285 \\
		\hline
		18 & argon & 39.948 & 1520.6 & 38.0644838289777 \\
		\hline
		19 & potassium & 39.0983 & 418.8 & 10.7114631582447 \\
		\hline
		20 & calcium & 40.078 & 589.8 & 14.7163032087429 \\
		\hline
		21 & scandium & 44.955907 & 633.1 & 14.0826877322262 \\
		\hline
		22 & titanium & 47.867 & 658.8 & 13.7631353542106 \\
		\hline
		23 & vanadium & 50.9415 & 650.9 & 12.7774015292051 \\
		\hline
		24 & chromium & 51.9961 & 652.9 & 12.556710984093 \\
		\hline
		25 & manganese & 54.938043 & 717.3 & 13.0565262399318 \\
		\hline
		26 & iron & 55.845 & 762.5 & 13.6538633718328 \\
		\hline
		27 & cobalt & 58.933194 & 760.4 & 12.9027454374864 \\
		\hline
		28 & nickel & 58.6934 & 737.1 & 12.5584818736008 \\
		\hline
		29 & copper & 63.546 & 745.5 & 11.7316589557171 \\
		\hline
		30 & zinc & 65.38 & 906.4 & 13.8635668400122 \\
		\hline
		31 & gallium & 69.723 & 578.8 & 8.3014213387261 \\
		\hline
		32 & germanium & 72.63 & 762 & 10.4915324246179 \\
		\hline
		33 & arsenic & 74.921595 & 947 & 12.6398803976344 \\
		\hline
		34 & selenium & 78.971 & 941 & 11.9157665472135 \\
		\hline
		35 & bromine & 79.904 & 1139.9 & 14.2658690428514 \\
		\hline
		36 & krypton & 83.798 & 1350.8 & 16.1197164610134 \\
		\hline
		37 & rubidium & 85.4678 & 403 & 4.71522608514552 \\
		\hline
		38 & strontium & 87.62 & 549.5 & 6.27139922392148 \\
		\hline
		39 & yttrium & 88.905838 & 600 & 6.7487131722441 \\
		\hline
		40 & zirconium & 91.224 & 640.1 & 7.01679382618609 \\
		\hline
		41 & niobium & 92.90637 & 652.1 & 7.01889439873714 \\
		\hline
		42 & molybdenum & 95.95 & 684.3 & 7.13183949973945 \\
		\hline
		43 & technetium & 98 & 702 & 7.16326530612245 \\
		\hline
		44 & ruthenium & 101.07 & 710.2 & 7.0268130998318 \\
		\hline
		45 & rhodium & 102.90549 & 719.7 & 6.99379595782499 \\
		\hline
		46 & palladium & 106.42 & 804.4 & 7.55872956211239 \\
		\hline
		47 & silver & 107.8682 & 731 & 6.77678871066728 \\
		\hline
		48 & cadmium & 112.414 & 867.8 & 7.7196790435355 \\
		\hline
		49 & indium & 114.818 & 558.3 & 4.8624780086746 \\
		\hline
		50 & tin & 118.71 & 708.6 & 5.96916856204195 \\
		\hline
		51 & antimony & 121.76 & 834 & 6.84954007884363 \\
		\hline
		52 & tellurium & 127.6 & 869.3 & 6.81269592476489 \\
		\hline
		53 & iodine & 126.90447 & 1008.4 & 7.9461346003021 \\
		\hline
		54 & xenon & 131.293 & 1170.4 & 8.91441280190109 \\
		\hline
		55 & caesium & 132.90545196 & 375.7 & 2.82682158225588 \\
		\hline
		56 & barium & 137.327 & 502.9 & 3.66206208538743 \\
		\hline
		57 & lanthanium & 138.90547 & 538.1 & 3.87385752339343 \\
		\hline
		58 & cerium & 140.116 & 534.4 & 3.81398270004853 \\
		\hline
		59 & praseodymium & 140.90766 & 527 & 3.74003797948245 \\
		\hline
		60 & neodymium & 144.242 & 533.1 & 3.69587221475021 \\
		\hline
	\end{tabular}
	\end{center}
	
	\begin{center}
	\begin{tabular}{|c|c|c|c|c|}
		\hline
		atomic number & element name & molar mass (gram/mol) & 1st ionization energy (kJ/mol) & specific ionization energy (MJ/kg) \\
		\hline
		61 & promethium & 145 & 540 & 3.72413793103448 \\
		\hline
		62 & samarium & 150.36 & 544.5 & 3.62130885873903 \\
		\hline
		63 & europium & 151.964 & 547.1 & 3.60019478297492 \\
		\hline
		64 & gadolinium & 157.25 & 593.4 & 3.77360890302067 \\
		\hline
		65 & terbium & 158.925354 & 565.8 & 3.56016196131927 \\
		\hline
		66 & dysprosium & 162.5 & 573 & 3.52615384615385 \\
		\hline
		67 & holmium & 164.930329 & 581 & 3.52269957577057 \\
		\hline
		68 & erbium & 167.259 & 589.3 & 3.52327826903186 \\
		\hline
		69 & thulium & 168.934219 & 596.7 & 3.53214407082321 \\
		\hline
		70 & ytterbium & 173.045 & 603.4 & 3.48695426045248 \\
		\hline
		71 & lutetium & 174.9668 & 523.5 & 2.99199619584973 \\
		\hline
		72 & hafnium & 178.486 & 658.5 & 3.68936499221227 \\
		\hline
		73 & tantalum & 180.94788 & 761 & 4.20563092532502 \\
		\hline
		74 & tungsten & 183.84 & 770 & 4.18842471714534 \\
		\hline
		75 & rhenium & 186.207 & 760 & 4.08147921399303 \\
		\hline
		76 & osmium & 190.23 & 840 & 4.41570730168743 \\
		\hline
		77 & iridium & 192.217 & 880 & 4.57815905981261 \\
		\hline
		78 & platinum & 195.084 & 870 & 4.45961739558344 \\
		\hline
		79 & gold & 196.96657 & 890.1 & 4.51904097228276 \\
		\hline
		80 & mercury & 200.59 & 1007.1 & 5.02068896754574 \\
		\hline
		81 & thallium & 204.38 & 589.4 & 2.88384382033467 \\
		\hline
		82 & lead & 207.2 & 715.6 & 3.45366795366795 \\
		\hline
		83 & bismuth & 208.9804 & 703 & 3.36395183471751 \\
		\hline
		84 & polonium & 209 & 812.1 & 3.88564593301435 \\
		\hline
		85 & astatine & 210 & 890 & 4.23809523809524 \\
		\hline
		86 & radon & 222 & 1037 & 4.67117117117117 \\
		\hline
		87 & francium & 223 & 380 & 1.70403587443946 \\
		\hline
		88 & radium & 226 & 509.3 & 2.25353982300885 \\
		\hline
		89 & actinium & 227 & 499 & 2.19823788546256 \\
		\hline
		90 & thorium & 232.0377 & 587 & 2.52976132757737 \\
		\hline
		91 & protactinium & 231.03588 & 568 & 2.45849259431046 \\
		\hline
		92 & uranium & 238.02891 & 597.6 & 2.51061940333214 \\
		\hline
		93 & neptunium & 237 & 604.5 & 2.55063291139241 \\
		\hline
		94 & plutonium & 244 & 584.7 & 2.39631147540984 \\
		\hline
		95 & americium & 243 & 578 & 2.37860082304527 \\
		\hline
		96 & curium & 247 & 581 & 2.35222672064777 \\
		\hline
		97 & berkelium & 247 & 601 & 2.4331983805668 \\
		\hline
		98 & californium & 251 & 608 & 2.42231075697211 \\
		\hline
		99 & einsteinium & 252 & 619 & 2.45634920634921 \\
		\hline
		100 & fermium & 257 & 627 & 2.43968871595331 \\
		\hline
		101 & mendelevium & 258 & 635 & 2.46124031007752 \\
		\hline
		102 & nobelium & 259 & 642 & 2.47876447876448 \\
		\hline
		103 & lawrencium & 266 & 470 & 1.76691729323308 \\
		\hline
		104 & ruthenfordium & 267 & 580 & 2.17228464419476 \\
		\hline
	\end{tabular}
	\end{center}
	
	\normalsize
	
	Average specific ionization energy: 32.3424893575117 MJ/kg.
	
	\medskip
	
	Taking a simple average of the specific ionization energy would correspond to a matter which is made of equal parts by weight of all these elements, what I refer to as "general matter".
	
	\medskip
	
	For a more accurate result, a weighted average should be taken, with the weights corresponding to the concentration by weight of the particular element in a given mix.
    
	\pagebreak
	
	\section{Specific ionization energy for meteor material}
	
	Here I am taking a weighted average of the specific ionization energies presented in the previous so that the final specific energy is more relevant to asteroid mining.
	
	\scriptsize
	
	\begin{center}
		\begin{tabular}{|c|c|c|c|}
			\hline
			atomic number & element name & normalized fraction of meteor material & adjusted specific ionization energy (MJ/kg) \\
			\hline
			1 & hydrogen & 0.0237970036865053 & 30.9738778141815 \\
			\hline
			2 & helium & 0 & 0 \\
			\hline
			3 & lithium & 1.69E-06 & 0.000126348716619381 \\
			\hline
			4 & beryllium & 2.88E-08 & 2.87E-06 \\
			\hline
			5 & boron & 1.59E-06 & 0.000117495412589677 \\
			\hline
			6 & carbon & 0.0148731273040658 & 1.34540444724565 \\
			\hline
			7 & nitrogen & 0.00138815854837948 & 0.138974422245487 \\
			\hline
			8 & oxygen & 0.396616728108421 & 32.571705672958 \\
			\hline
			9 & fluorine & 8.63E-05 & 0.00763274762371741 \\
			\hline
			10 & neon & 0 & 0 \\
			\hline
			11 & sodium & 0.0054534800114908 & 0.117610375152801 \\
			\hline
			12 & magnesium & 0.118985018432526 & 3.6114070396081 \\
			\hline
			13 & aluminium & 0.00902303056446659 & 0.193124649667102 \\
			\hline
			14 & silicon & 0.138815854837948 & 3.88743705999807 \\
			\hline
			15 & phosphorus & 0.00109069600229816 & 0.0356290661494892 \\
			\hline
			16 & sulfur & 0.0396616728108422 & 1.23661285532495 \\
			\hline
			17 & chlorine & 0.00036687047350029 & 0.0129486131577874 \\
			\hline
			18 & argon & 0 & 0 \\
			\hline
			19 & potassium & 0.000694079274189738 & 0.00743460457438462 \\
			\hline
			20 & calcium & 0.0109069600229816 & 0.160510130783835 \\
			\hline
			21 & scandium & 6.35E-06 & 8.94E-05 \\
			\hline
			22 & titanium & 0.000535432582946369 & 0.00736923111214548 \\
			\hline
			23 & vanadium & 6.05E-05 & 0.000772829006206731 \\
			\hline
			24 & chromium & 0.00297462546081316 & 0.0373515121973554 \\
			\hline
			25 & manganese & 0.00267716291473185 & 0.0349544478447686 \\
			\hline
			26 & iron & 0.218139200459632 & 2.97844283911665 \\
			\hline
			27 & cobalt & 0.000585009673959922 & 0.00754823090157178 \\
			\hline
			28 & nickel & 0.0128900436635237 & 0.161879379698285 \\
			\hline
			29 & copper & 0.000109069600229816 & 0.00127956735233261 \\
			\hline
			30 & zinc & 0.00017847752764879 & 0.00247433513399913 \\
			\hline
			31 & gallium & 7.54E-06 & 6.26E-05 \\
			\hline
			32 & germanium & 2.08E-05 & 0.000218458656312507 \\
			\hline
			33 & arsenic & 1.78E-06 & 2.26E-05 \\
			\hline
			34 & selenium & 1.29E-05 & 0.000153594751077937 \\
			\hline
			35 & bromine & 1.19E-06 & 1.70E-05 \\
			\hline
			36 & krypton & 0 & 0 \\
			\hline
			37 & rubidium & 3.17E-06 & 1.50E-05 \\
			\hline
			38 & strontium & 8.63E-06 & 5.41E-05 \\
			\hline
			39 & yttrium & 1.88E-06 & 1.27E-05 \\
			\hline
			40 & zirconium & 6.54E-06 & 4.59E-05 \\
			\hline
			41 & niobium & 1.88E-07 & 1.32E-06 \\
			\hline
			42 & molybdenum & 1.19E-06 & 8.49E-06 \\
			\hline
			43 & technetium & 0 & 0 \\
			\hline
			44 & ruthenium & 8.03E-07 & 5.64E-06 \\
			\hline
			45 & rhodium & 1.78E-07 & 1.25E-06 \\
			\hline
			46 & palladium & 6.54E-07 & 4.95E-06 \\
			\hline
			47 & silver & 1.39E-07 & 9.41E-07 \\
			\hline
			48 & cadmium & 4.36E-07 & 3.37E-06 \\
			\hline
			49 & indium & 4.36E-08 & 2.12E-07 \\
			\hline
			50 & tin & 1.19E-06 & 7.10E-06 \\
			\hline
			51 & antimony & 1.19E-07 & 8.15E-07 \\
			\hline
			52 & tellurium & 2.08E-06 & 1.42E-05 \\
			\hline
			53 & iodine & 2.48E-07 & 1.97E-06 \\
			\hline
			54 & xenon & 0 & 0 \\
			\hline
			55 & caesium & 1.39E-07 & 3.92E-07 \\
			\hline
			56 & barium & 2.68E-06 & 9.80E-06 \\
			\hline
			57 & lanthanium & 2.78E-07 & 1.08E-06 \\
			\hline
			58 & cerium & 7.44E-07 & 2.84E-06 \\
			\hline
			59 & praseodymium & 9.72E-08 & 3.63E-07 \\
			\hline
			60 & neodymium & 4.96E-07 & 1.83E-06 \\
			\hline
		\end{tabular}
	\end{center}
	
	\begin{center}
		\begin{tabular}{|c|c|c|c|}
			\hline
			atomic number & element name & normalized fraction of meteor material & adjusted specific ionization energy (MJ/kg) \\
			\hline
			61 & promethium & 0 & 0 \\
			\hline
			62 & samarium & 1.69E-07 & 6.10E-07 \\
			\hline
			63 & europium & 5.85E-08 & 2.11E-07 \\
			\hline
			64 & gadolinium & 2.28E-07 & 8.61E-07 \\
			\hline
			65 & terbium & 3.87E-08 & 1.38E-07 \\
			\hline
			66 & dysprosium & 2.68E-07 & 9.44E-07 \\
			\hline
			67 & holmium & 5.85E-08 & 2.06E-07 \\
			\hline
			68 & erbium & 1.78E-07 & 6.29E-07 \\
			\hline
			69 & thulium & 2.88E-08 & 1.02E-07 \\
			\hline
			70 & ytterbium & 1.78E-07 & 6.22E-07 \\
			\hline
			71 & lutetium & 2.88E-08 & 8.60E-08 \\
			\hline
			72 & hafnium & 1.69E-07 & 6.22E-07 \\
			\hline
			73 & tantalum & 1.98E-08 & 8.34E-08 \\
			\hline
			74 & tungsten & 1.19E-07 & 4.98E-07 \\
			\hline
			75 & rhenium & 4.86E-08 & 1.98E-07 \\
			\hline
			76 & osmium & 6.54E-07 & 2.89E-06 \\
			\hline
			77 & iridium & 5.35E-07 & 2.45E-06 \\
			\hline
			78 & platinum & 9.72E-07 & 4.33E-06 \\
			\hline
			79 & gold & 1.69E-07 & 7.62E-07 \\
			\hline
			80 & mercury & 2.48E-07 & 1.24E-06 \\
			\hline
			81 & thallium & 7.73E-08 & 2.23E-07 \\
			\hline
			82 & lead & 1.39E-06 & 4.79E-06 \\
			\hline
			83 & bismuth & 6.84E-08 & 2.30E-07 \\
			\hline
			84 & polonium & 0 & 0 \\
			\hline
			85 & astatine & 0 & 0 \\
			\hline
			86 & radon & 0 & 0 \\
			\hline
			87 & francium & 0 & 0 \\
			\hline
			88 & radium & 0 & 0 \\
			\hline
			89 & actinium & 0 & 0 \\
			\hline
			90 & thorium & 3.87E-08 & 9.78E-08 \\
			\hline
			91 & protactinium & 0 & 0 \\
			\hline
			92 & uranium & 9.72E-09 & 2.44E-08 \\
			\hline
			93 & neptunium & 0 & 0 \\
			\hline
			94 & plutonium & 0 & 0 \\
			\hline
			95 & americium & 0 & 0 \\
			\hline
			96 & curium & 0 & 0 \\
			\hline
			97 & berkelium & 0 & 0 \\
			\hline
			98 & californium & 0 & 0 \\
			\hline
			99 & einsteinium & 0 & 0 \\
			\hline
			100 & fermium & 0 & 0 \\
			\hline
			101 & mendelevium & 0 & 0 \\
			\hline
			102 & nobelium & 0 & 0 \\
			\hline
			103 & lawrencium & 0 & 0 \\
			\hline
			104 & ruthenfordium & 0 & 0 \\
			\hline
		\end{tabular}
	\end{center}
	
	\normalsize
	
	In the table above the normalized fraction of meteor material is just the mass fraction of the given element in a typical meteor, adjusted so that the sum of all values in that column is 1.
	
	\medskip
	
	The adjusted specific ionization energy is the ionization energy of that material multiplied by the mass fraction.
	
	\medskip
	
	The sum of the adjusted specific ionization energies should in that case be the energy needed to ionize meteor material per kilogram.
	It adds up to 77.5334091946941 MJ/kg.
	
	\medskip
	
	Playing with the libre calc sheet used to calculate this revealed that this final sum is disproportionately inflated by the almost 2.4\% of hydrogen included in the meteor material.
	If the hydrogen were removed and everything else adjusted up to fill in the gap, it would only be a bit over 44MJ/kg.
	
	\pagebreak
	
	\section{Specific ionization energy for electrolytic Kamacite}
	
	This is a particular special case for a material I use a lot.
	See one of my other documents "Felicity tech tree: low temperature electrolysis of mixed metal oxides" for context to this.
	
	\tiny
	
	\begin{center}
		\begin{tabular}{|c|c|c|c|c|c|}
			\hline
			atomic number & element name & specific ionization energy (MJ/kg) & mass in LTE batch & normalized kamacite mass fraction & adjusted ionization energy (MJ/kg) \\
			\hline
			26 & iron & 13.6538633718328 & 22 & 0.940433819556142 & 12.8405548824704 \\
			\hline
			27 & cobalt & 12.9027454374864 & 0.059 & 0.00252207251608238 & 0.0325416596498918 \\
			\hline
			28 & nickel & 12.5584818736008 & 1.3 & 0.0555710893374084 & 0.697888518140094 \\
			\hline
			29 & copper & 11.7316589557171 & 0.011 & 0.000470216909778071 & 0.00551642442072752 \\
			\hline
			30 & zinc & 13.8635668400122 & 0.018 & 0.000769445852364116 & 0.0106672640040201 \\
			\hline
			31 & gallium & 8.3014213387261 & 0.00076 & 3.24877137664849E-05 & 0.000269694200307523 \\
			\hline
			32 & germanium & 10.4915324246179 & 0.0021 & 8.97686827758135E-05 & 0.000941811046057686 \\
			\hline
			33 & arsenic & 12.6398803976344 & 0.00018 & 7.69445852364116E-06 & 9.72570354633828E-05 \\
			\hline
			34 & selenium & 11.9157665472135 & 0.0013 & 5.55710893374084E-05 & 0.000662172127318904 \\
			\hline
			42 & molybdenum & 7.13183949973945 & 0.00012 & 5.12963901576077E-06 & 3.65837621520073E-05 \\
			\hline
			44 & ruthenium & 7.0268130998318 & 0.000081 & 3.46250633563852E-06 & 2.43303848775154E-05 \\
			\hline
			45 & rhodium & 6.99379595782499 & 0.000016 & 6.83951868768103E-07 & 4.78341981513721E-06 \\
			\hline
			46 & palladium & 7.55872956211239 & 0.000066 & 2.82130145866843E-06 & 2.13254547392678E-05 \\
			\hline
			47 & silver & 6.77678871066728 & 0.000014 & 5.9845788517209E-07 & 4.05562264004404E-06 \\
			\hline
			48 & cadmium & 7.7196790435355 & 0.000044 & 1.88086763911228E-06 & 1.45196944973192E-05 \\
			\hline
			49 & indium & 4.8624780086746 & 0.0000044 & 1.88086763911228E-07 & 9.14567753241119E-07 \\
			\hline
			50 & tin & 5.96916856204195 & 0.00012 & 5.12963901576077E-06 & 3.0619679947503E-05 \\
			\hline
			51 & antimony & 6.84954007884363 & 0.000012 & 5.12963901576077E-07 & 3.51356680284534E-06 \\
			\hline
			52 & tellurium & 6.81269592476489 & 0.00021 & 8.97686827758135E-06 & 6.11566739318297E-05 \\
			\hline
			74 & tungsten & 4.18842471714534 & 0.000012 & 5.12963901576077E-07 & 2.14851068436455E-06 \\
			\hline
			75 & rhenium & 4.08147921399303 & 0.0000049 & 2.09460259810232E-07 & 8.5490769657304E-07 \\
			\hline
			76 & osmium & 4.41570730168743 & 0.000066 & 2.82130145866843E-06 & 1.24580414513036E-05 \\
			\hline
			77 & iridium & 4.57815905981261 & 0.000054 & 2.30833755709235E-06 & 1.0567936500108E-05 \\
			\hline
			78 & platinum & 4.45961739558344 & 0.000098 & 4.18920519620463E-06 & 1.86822523666627E-05 \\
			\hline
			79 & gold & 4.51904097228276 & 0.000017 & 7.2669886056611E-07 & 3.28398192540945E-06 \\
			\hline
			80 & mercury & 5.02068896754574 & 0.000025 & 1.06867479495016E-06 & 5.36548375290048E-06 \\
			\hline
			81 & thallium & 2.88384382033467 & 0.0000078 & 3.3342653602445E-07 & 9.61550055449706E-07 \\
			\hline
			82 & lead & 3.45366795366795 & 0.00014 & 5.9845788517209E-06 & 2.06687481963874E-05 \\
			\hline
			83 & bismuth & 3.36395183471751 & 0.0000069 & 2.94954243406244E-07 & 9.92211868264151E-07 \\
			\hline
		\end{tabular}
	\end{center}
	
	\normalsize
	
	The sum of adjusted ionization energy for Kamacite turns out to be 13.589417469546 MJ/kg.
	
	\pagebreak
	
	\section{Calutron efficiency}
	
    The total efficiency of the calutron here means the ratio of energy per kilogram that is required to ionize the material against the energy that is actually used by the caluton per kilogram of material refined.
    
    \medskip
    
    It should be a product of all the following factors multiplied together:
    
    \begin{itemize}
    	\item Efficiency of the excimer laser used for initial ionization of the material: between 2\% and 0.2\% (\href{https://spie.org/publications/spie-publication-resources/optipedia-free-optics-information/fg12_p87-88_excimer_lasers}{spie.org})
    	
    	\item Efficiency of the ion accelerator: unknown, for now assumed 50\%.
    	
    	\item Efficiency multipllier of by the deflector magnet: unknown, for now assumed 50\%.
    \end{itemize}
    
    In the best case this gives a final efficiency of 0.5\%, and in the worst case 0.05\%.
    
    $$ \eta_{optimistic} = 0.02 \cdot 0.5 \cdot 0.5 = 0.005 $$
    
    $$ \eta_{pessimistic} = 0.002 \cdot 0.5 \cdot 0.5 = 0.005 $$
    
    Correspondingly the specific energy of refining general matter (composed in equal parts by weight of the first 104 elements of the periodic table) would be approximately between 6.4685 GJ/kg and 64.685 GJ/kg.
    
    $$ \epsilon_{optimistic} = \frac{32.3424893575117 MJ/kg}{0.005} = 6.468497872 GJ/kg $$
    
    $$ \epsilon_{pessimistic} = \frac{32.3424893575117 MJ/kg}{0.0005} = 64.68497872 GJ/kg $$
    
    For meteor material the energy per kilogram to be refined this way is between 15.5GJ/kg and 155GJ/kg .
    
    $$ \epsilon_{optimistic} = \frac{77.5334091946941 MJ/kg}{0.005} = 15.50668184 GJ/kg $$
    
    $$ \epsilon_{pessimistic} = \frac{77.5334091946941 MJ/kg}{0.0005} = 155.0668184 GJ/kg $$
    
    For electrolytic kamacite the energy per kilogram to be refined this way is between 2.71788GJ/kg and 27.1788GJ/kg .
    
    $$ \epsilon_{optimistic} = \frac{13.589417469546 MJ/kg}{0.005} = 2.717883494 GJ/kg $$
    
    $$ \epsilon_{pessimistic} = \frac{13.589417469546 MJ/kg}{0.0005} = 27.17883494 GJ/kg $$
	
\end{document}